\section{Baccalauréat série C — Session 2023}
\noindent\begin{minipage}{0.5\linewidth}\footnotesize Office du Baccalauréat du Cameroun\end{minipage}%
\hfill\begin{minipage}{0.45\linewidth}\footnotesize\raggedleft Examen : \textbf{Baccalauréat} · Série : \textbf{C}\\
Épreuve : \textbf{Mathématiques} · Durée : \textbf{4\,h} · Coef. : \textbf{7}\end{minipage}
\par\smallskip{\color{onbo}\rule{\linewidth}{1pt}}

\subsection*{Partie A — Évaluation des ressources \hfill (13,25 points)}

\noindent\textbf{Exercice 1.} \hfill\textit{(4,5 points)}\\
On considère les fonctions numériques $f$ et $h$ de la variable réelle $x$ définies sur
$\mathcal D=\,]1\,;\,+\infty[$ par
\[
f(x)=x-2+\ln\sqrt{x-1}
\qquad\text{et}\qquad
h(x)=\frac{2x-2-\ln(x-1)}{2\sqrt{x-1}}.
\]
Le plan est rapporté à un repère orthonormé $(O\,;\,\vec\imath,\vec\jmath)$.
\begin{enumerate}[label=\arabic*)]
\item Dresser le tableau de variations de $f$ sur $\mathcal D$. \hfill\textit{(0,75 pt)}
\item Montrer que le réel $2$ est l'unique solution de l'équation $f(x)=0$. \hfill\textit{(0,5 pt)}
\item En déduire, suivant les valeurs de $x$, le signe de $f(x)$. \hfill\textit{(0,5 pt)}
\item Montrer que pour tout $x\in\,]1\,;\,+\infty[$,\ $\displaystyle h'(x)=\frac{f(x)}{2(x-1)\sqrt{x-1}}$,
puis en déduire les variations de $h$. \hfill\textit{(0,75 pt)}
\item On considère la suite $(u_n)_n$ définie par
$\displaystyle u_n=\sum_{j=0}^{n}\frac1n\,h\!\left(2+\frac jn\right)$ et on pose
$\displaystyle I=\int_2^3 h(x)\dd x$.
  \begin{enumerate}[label=\alph*)]
  \item Calculer $\displaystyle\int_2^3 \frac{\ln(x-1)}{2\sqrt{x-1}}\dd x$ à l'aide d'une intégration
  par parties et en déduire la valeur de $I$. \hfill\textit{(0,75 pt)}
  \item Soient $n\in\N^*$ et $j$ un entier naturel tel que $0\le j\le n-1$. En utilisant les
  variations de $h$ sur $[2\,;\,+\infty[$, démontrer que
  \[
  \frac1n\,h\!\left(2+\frac jn\right)\le\int_{2+\frac jn}^{2+\frac{j+1}{n}}h(x)\dd x\le\frac1n\,h\!\left(2+\frac{j+1}{n}\right).
  \]
  \hfill\textit{(0,5 pt)}
  \item Déduire de la question précédente que :
  $\displaystyle u_n-\frac{h(3)}{n}\le I\le u_n-\frac{h(2)}{n}$. \hfill\textit{(0,5 pt)}
  \item Calculer la limite de la suite $(u_n)_{n\in\N}$. \hfill\textit{(0,25 pt)}
  \end{enumerate}
\end{enumerate}

\smallskip
\noindent\textbf{Exercice 2.} \hfill\textit{(4,25 points)}\\
Le plan est muni d'un repère orthonormé direct $(O\,;\,\vec\imath,\vec\jmath)$. On considère
l'équation
\[
(E)\,:\ z^2+\bigl(-3\cos\alpha-1+\mathrm i(3-5\sin\alpha)\bigr)z+5\sin\alpha-2+\mathrm i(-3\cos\alpha-1)=0
\]
d'inconnue complexe $z$, où $\alpha$ est un nombre réel.
\begin{enumerate}[label=\arabic*)]
\item Montrer que $-\mathrm i$ est une solution de $(E)$. \hfill\textit{(0,25 pt)}
\item En déduire l'autre solution. \hfill\textit{(0,5 pt)}
\item Montrer que l'ensemble des points $A_\alpha$ d'affixe
$z_\alpha=3\cos\alpha+1+\mathrm i(-2+5\sin\alpha)$ lorsque $\alpha$ décrit $\R$ est la conique
$(\epsilon)$ d'équation : $25x^2+9y^2-50x+36y-164=0$. \hfill\textit{(0,5 pt)}
\item Soit $\Omega(1\,;\,-2)$ un point du plan.
  \begin{enumerate}[label=\alph*)]
  \item Déterminer une équation de $(\epsilon)$ dans le repère $(\Omega\,;\,\vec\imath,\vec\jmath)$. \hfill\textit{(0,5 pt)}
  \item En déduire la nature exacte de $(\epsilon)$ ; préciser son excentricité et les
  coordonnées de ses sommets dans le repère $(\Omega\,;\,\vec\imath,\vec\jmath)$. \hfill\textit{(1 pt)}
  \item Construire $(\epsilon)$ dans le repère $(\Omega\,;\,\vec\imath,\vec\jmath)$. \hfill\textit{(0,75 pt)}
  \end{enumerate}
\item Soient $B$ et $C$ deux points d'affixes respectives $-\mathrm i$ et $3$.
  \begin{enumerate}[label=\alph*)]
  \item Déterminer l'écriture complexe de la similitude directe $S$ de centre $\Omega$ telle que $S(C)=B$. \hfill\textit{(0,5 pt)}
  \item En déduire l'angle de $S$. \hfill\textit{(0,25 pt)}
  \end{enumerate}
\end{enumerate}

\smallskip
\noindent\textbf{Exercice 3.} \hfill\textit{(4,5 points)}\\
L'espace est muni d'un repère orthonormé direct $(O\,;\,\vec\imath,\vec\jmath,\vec k)$.
Soient $A(-1\,;\,-1\,;\,0)$, $B(0\,;\,0\,;\,2)$ et $C(-1\,;\,1\,;\,2)$ trois points de l'espace.
\begin{enumerate}[label=\arabic*)]
\item Montrer que les points $A$, $B$ et $C$ définissent un plan. \hfill\textit{(0,5 pt)}
\item Déterminer une équation cartésienne de ce plan. \hfill\textit{(0,5 pt)}
\item Soit $(P)$ le plan d'équation $x+y-z+2=0$. Déterminer l'expression analytique de la
réflexion $f$ de plan $(P)$. \hfill\textit{(0,75 pt)}
\item Soit $g$ la transformation de l'espace d'expression analytique :
\[
\begin{cases}
x'=\dfrac13(-x+2y-2z+4)\\[4pt]
y'=\dfrac13(2x-y-2z+4)\\[4pt]
z'=\dfrac13(-2x-2y-z+8)
\end{cases}
\]
  \begin{enumerate}[label=\alph*)]
  \item Montrer que l'ensemble $(D)$ des points invariants par $g$ est la droite passant par
  $B$ dont un vecteur directeur est $\vec u(-1\,;\,-1\,;\,1)$. \hfill\textit{(0,75 pt)}
  \item Soient $M$ et $M'$ deux points de l'espace tels que $g(M)=M'$.
    \begin{enumerate}[label=\roman*)]
    \item Montrer que $\overrightarrow{MM'}$ est un vecteur normal à la droite $(D)$. \hfill\textit{(0,5 pt)}
    \item Montrer que le milieu du segment $[MM']$ appartient à $(D)$. \hfill\textit{(0,5 pt)}
    \end{enumerate}
  \item En déduire que $g$ est un demi-tour. \hfill\textit{(0,25 pt)}
  \end{enumerate}
\item
  \begin{enumerate}[label=\alph*)]
  \item Montrer que $(P)\perp(D)$. \hfill\textit{(0,25 pt)}
  \item En déduire que $f\circ g$ est une symétrie centrale dont on précisera le centre. \hfill\textit{(0,5 pt)}
  \end{enumerate}
\end{enumerate}

\subsection*{Partie B — Évaluation des compétences \hfill (6,75 points)}

\noindent\textbf{Situation.}\par\smallskip
\noindent\textit{KEMO vend des marchandises qu'il pèse sur une balance très contestée par sa
clientèle ces derniers jours. Pour le convaincre, il envisage d'acquérir une nouvelle balance
constituée d'un ressort que l'on suspend verticalement, pouvant s'étirer ou s'allonger d'au plus
de \SI{7}{\centi\metre}, et voudrait investir \num{130000}~FCFA dans la publicité. Il lui faut
cependant convaincre ses deux fournisseurs qui lui rendent visite à des fréquences différentes.
Pour cela, il a besoin de connaître : la masse maximale que peut peser la balance sollicitée, le
chiffre d'affaires que cette publicité pourrait lui permettre de réaliser et la prochaine date de
coïncidence de la visite des deux fournisseurs.}

\noindent\textbf{Données.}
\begin{itemize}
\item[\textit{a)}] \textit{Du ressort de la balance.} Une étude expérimentale montre que le ressort
est indéformable et s'allonge de \SI{2}{\centi\metre} lorsqu'on y accroche une masse de
\SI{4}{\kilo\gram}. Par ailleurs, lorsqu'on l'étire de sa position d'équilibre et qu'on l'abandonne
sans vitesse initiale, son élongation $x(t)$ vérifie l'équation
$x''(t)+\dfrac k4\,x(t)=0$, où $k$ est la constante de raideur du ressort. De plus, on a l'égalité
$mg=k\,\Delta\ell_0$, où $m$ est la masse du corps accroché au ressort et $\Delta\ell_0$
l'allongement au repos. Après une minute, le centre de gravité du solide repasse pour la première
fois au point initial. On prend $g=\SI{9,5}{\newton\per\kilo\gram}$ et $\pi=3{,}14$.
\item[\textit{b)}] \textit{Chiffre d'affaires en fonction des frais de publicité.} Le tableau
ci-dessous donne les dépenses en publicité (en dizaines de milliers) et le chiffre d'affaires
correspondant (en dizaines de millions), sur les dix dernières années. On admettra que le chiffre
d'affaires suit un ajustement linéaire par rapport aux frais de publicité.
\par\smallskip
\begin{center}\footnotesize
\begin{tabular}{|l|*{10}{c|}}
\hline
Frais de publicité $(x_i)$ & 6 & 6,5 & 6,8 & 7 & 7,8 & 9 & 10,5 & 11 & 11,3 & 11\\
\hline
Chiffre d'affaires $(y_i)$ & 220 & 229 & 225 & 237 & 235 & 247 & 250 & 268 & 258 & 264\\
\hline
\end{tabular}
\end{center}
\item[\textit{c)}] Le premier fournisseur lui rend visite tous les $21$ jours et le
\textbf{20 décembre 2020} il était au marché. Quant au second fournisseur, il lui rend visite tous
les $16$ jours et était au marché le \textbf{27 décembre 2020}.
\end{itemize}

\noindent\textbf{Tâches.}
\begin{enumerate}[label=\textbf{Tâche \arabic*.},leftmargin=2.4em]
\item Déterminer la masse maximale que cette balance peut peser. \hfill\textit{(2,25 pts)}
\item Estimer le chiffre d'affaires qu'il pourra espérer des frais de publicité investis. \hfill\textit{(2,25 pts)}
\item Donner la date de la prochaine coïncidence des deux fournisseurs. \hfill\textit{(2,25 pts)}
\end{enumerate}

% ════════════════════════════════════════════════════
\corrige{de l'annale 2023}

\noindent\textbf{Partie A — Exercice 1.}

\noindent\textbf{1)} \emph{Limites aux bornes de $\mathcal D=\,]1\,;\,+\infty[$.}
Quand $x\to1^+$, $\sqrt{x-1}\to0^+$ donc $\ln\sqrt{x-1}\to-\infty$ et $x-2\to-1$, d'où
$\displaystyle\lim_{x\to1^+}f(x)=-\infty$.
Quand $x\to+\infty$, en posant $X=x-1$,
\[
f(x)=X-1+\ln\sqrt X=\ln\!\left(\mathrm e^{X}\sqrt X\right)-1\longrightarrow+\infty.
\]
Donc $\displaystyle\lim_{x\to+\infty}f(x)=+\infty$.

\noindent\emph{Dérivée.} $f$ est dérivable sur $\mathcal D$ et
\[
f'(x)=(x-2)'+(\ln\sqrt{x-1})'=1+\frac{(\sqrt{x-1})'}{\sqrt{x-1}}=1+\frac{\frac{1}{2\sqrt{x-1}}}{\sqrt{x-1}}=1+\frac1{2(x-1)}.
\]
Pour $x>1$, on a $x-1>0$, donc $1+\dfrac1{2(x-1)}>0$ : $f'(x)>0$. La fonction $f$ est
strictement croissante sur $\mathcal D$.
\begin{center}
\begin{tabular}{c|ccc}
$x$ & $1$ & & $+\infty$\\\hline
$f'(x)$ & & $+$ & \\\hline
$f(x)$ & $-\infty$ & $\nearrow$ & $+\infty$\\
\end{tabular}
\end{center}

\noindent\textbf{2)} $f$ est continue et strictement croissante sur $]1\,;\,+\infty[$ ; elle
réalise donc une bijection de $]1\,;\,+\infty[$ sur $\R$. Or
$f(2)=2-2+\ln\sqrt{2-1}=\ln1=0$. Donc $2$ est l'unique solution de l'équation $f(x)=0$.

\noindent\textbf{3)} $f$ étant strictement croissante et s'annulant en $2$ :
\begin{itemize}
\item si $x\in\,]1\,;\,2[$, alors $f(x)<0$ ;
\item si $x\in\,]2\,;\,+\infty[$, alors $f(x)>0$.
\end{itemize}

\noindent\textbf{4)} $h$ est dérivable sur $]1\,;\,+\infty[$. En dérivant le quotient
$h=\dfrac{2x-2-\ln(x-1)}{2\sqrt{x-1}}$, avec
$(2x-2-\ln(x-1))'=2-\dfrac1{x-1}=\dfrac{2x-3}{x-1}$ et $(2\sqrt{x-1})'=\dfrac1{\sqrt{x-1}}$ :
\[
h'(x)=\frac{\dfrac{2x-3}{x-1}\cdot 2\sqrt{x-1}-\bigl(2x-2-\ln(x-1)\bigr)\cdot\dfrac1{\sqrt{x-1}}}{4(x-1)}
=\frac{\dfrac{2(2x-3)}{\sqrt{x-1}}-\dfrac{2x-2-\ln(x-1)}{\sqrt{x-1}}}{4(x-1)}.
\]
Le numérateur vaut $\dfrac{4x-6-(2x-2)+\ln(x-1)}{\sqrt{x-1}}=\dfrac{2x-4+\ln(x-1)}{\sqrt{x-1}}$, donc
\[
h'(x)=\frac{2x-4+\ln(x-1)}{4(x-1)\sqrt{x-1}}
=\frac{2x-4+2\ln\sqrt{x-1}}{4(x-1)\sqrt{x-1}}
=\frac{2\bigl(x-2+\ln\sqrt{x-1}\bigr)}{4(x-1)\sqrt{x-1}}.
\]
Comme $f(x)=x-2+\ln\sqrt{x-1}$, on obtient bien
\[
\boxed{\,h'(x)=\frac{f(x)}{2(x-1)\sqrt{x-1}}\,}.
\]
Pour $x>1$, $2(x-1)\sqrt{x-1}>0$, donc $h'(x)$ a le signe de $f(x)$ (question 3). Par conséquent :
\begin{itemize}
\item $h$ est strictement décroissante sur $]1\,;\,2[$ ;
\item $h$ est strictement croissante sur $]2\,;\,+\infty[$,
\end{itemize}
avec un minimum en $x=2$.

\noindent\textbf{5) a)} \emph{Intégration par parties} de
$\displaystyle J=\int_2^3\frac{\ln(x-1)}{2\sqrt{x-1}}\dd x$, avec
$u(x)=\ln(x-1)$, $u'(x)=\dfrac1{x-1}$ et $v'(x)=\dfrac1{2\sqrt{x-1}}$, $v(x)=\sqrt{x-1}$ :
\[
J=\Bigl[\sqrt{x-1}\,\ln(x-1)\Bigr]_2^3-\int_2^3\frac{1}{x-1}\sqrt{x-1}\dd x
=\Bigl[\sqrt{x-1}\,\ln(x-1)\Bigr]_2^3-\int_2^3\frac{1}{\sqrt{x-1}}\dd x.
\]
Comme $\displaystyle\int_2^3\frac{\dd x}{\sqrt{x-1}}=\bigl[2\sqrt{x-1}\bigr]_2^3$, on a
\[
J=\Bigl[(\ln(x-1)-2)\sqrt{x-1}\Bigr]_2^3=(\ln2-2)\sqrt2-(\ln1-2)\cdot1=(\ln2-2)\sqrt2+2.
\]
\emph{Valeur de $I$.}
\[
I=\int_2^3\frac{2x-2-\ln(x-1)}{2\sqrt{x-1}}\dd x
=\int_2^3\frac{2(x-1)}{2\sqrt{x-1}}\dd x-\int_2^3\frac{\ln(x-1)}{2\sqrt{x-1}}\dd x
=\int_2^3\sqrt{x-1}\dd x-J.
\]
Or $\displaystyle\int_2^3\sqrt{x-1}\dd x=\Bigl[\tfrac23(x-1)^{3/2}\Bigr]_2^3=\tfrac23\bigl(2\sqrt2-1\bigr)=\frac{4\sqrt2}{3}-\frac23$. Ainsi
\[
I=\frac{4\sqrt2}{3}-\frac23-\bigl[(\ln2-2)\sqrt2+2\bigr]
=\frac{4\sqrt2}{3}-\frac23-\sqrt2\ln2+2\sqrt2-2
=\boxed{\,\frac{10\sqrt2-8}{3}-\sqrt2\ln2\,}.
\]

\noindent\textbf{5) b)} Soit $x\in\Bigl[2+\dfrac jn\,;\,2+\dfrac{j+1}{n}\Bigr]$. Comme $h$ est
croissante sur $[2\,;\,+\infty[$,
\[
h\!\left(2+\frac jn\right)\le h(x)\le h\!\left(2+\frac{j+1}{n}\right).
\]
En intégrant sur l'intervalle de longueur $\dfrac1n$ :
\[
\frac1n\,h\!\left(2+\frac jn\right)\le\int_{2+\frac jn}^{2+\frac{j+1}{n}}h(x)\dd x\le\frac1n\,h\!\left(2+\frac{j+1}{n}\right).
\]

\noindent\textbf{5) c)} En sommant l'encadrement précédent pour $j$ de $0$ à $n-1$ et grâce à la
relation de Chasles ($\sum_{j=0}^{n-1}\int_{2+\frac jn}^{2+\frac{j+1}{n}}h=\int_2^3 h=I$) :
\[
\sum_{j=0}^{n-1}\frac1n\,h\!\left(2+\frac jn\right)\le I\le\sum_{j=0}^{n-1}\frac1n\,h\!\left(2+\frac{j+1}{n}\right).
\]
À gauche, $\displaystyle\sum_{j=0}^{n-1}\frac1n h\!\left(2+\frac jn\right)=u_n-\frac1n h(3)$
(en complétant la somme jusqu'à $j=n$, qui donne le terme $\tfrac1n h(2+\tfrac nn)=\tfrac1n h(3)$).
À droite, le changement d'indice $k=j+1$ donne
$\displaystyle\sum_{j=0}^{n-1}\frac1n h\!\left(2+\frac{j+1}{n}\right)=u_n-\frac1n h(2)$. D'où
\[
\boxed{\,u_n-\frac{h(3)}{n}\le I\le u_n-\frac{h(2)}{n}\,}.
\]

\noindent\textbf{5) d)} De l'encadrement précédent on tire
$\dfrac{h(2)}{n}\le u_n-I\le\dfrac{h(3)}{n}$. Comme
$\displaystyle\lim_{n\to+\infty}\frac{h(2)}{n}=\lim_{n\to+\infty}\frac{h(3)}{n}=0$, le théorème des
gendarmes donne $\displaystyle\lim_{n\to+\infty}u_n=I$, soit
\[
\boxed{\,\lim_{n\to+\infty}u_n=I=\frac{10\sqrt2-8}{3}-\sqrt2\ln2\,}.
\]

\bigskip
\noindent\textbf{Exercice 2.}

\noindent\textbf{1)} En remplaçant $z$ par $-\mathrm i$ dans le membre de gauche de $(E)$ :
\[
(-\mathrm i)^2+\bigl(-3\cos\alpha-1+\mathrm i(3-5\sin\alpha)\bigr)(-\mathrm i)+5\sin\alpha-2+\mathrm i(-3\cos\alpha-1)
\]
\[
=-1+3\mathrm i\cos\alpha+\mathrm i+(3-5\sin\alpha)+5\sin\alpha-2-3\mathrm i\cos\alpha-\mathrm i=0.
\]
Donc $-\mathrm i$ est solution de $(E)$.

\noindent\textbf{2)} Si $z_1,z_2$ sont les racines de $z^2+bz+c=0$, alors $z_1+z_2=-b$. Ici
$z_1=-\mathrm i$ et $b=-3\cos\alpha-1+\mathrm i(3-5\sin\alpha)$, d'où
\[
z_2=-b-z_1=3\cos\alpha+1-\mathrm i(3-5\sin\alpha)+\mathrm i=3\cos\alpha+1+\mathrm i(-2+5\sin\alpha).
\]

\noindent\textbf{3)} Posons $z_\alpha=x+\mathrm iy$. Alors $x=3\cos\alpha+1$ et $y=-2+5\sin\alpha$,
soit $\cos\alpha=\dfrac{x-1}{3}$ et $\sin\alpha=\dfrac{y+2}{5}$. La relation
$\cos^2\alpha+\sin^2\alpha=1$ donne
\[
\left(\frac{x-1}{3}\right)^2+\left(\frac{y+2}{5}\right)^2=1
\;\Longrightarrow\;
\frac{25(x-1)^2+9(y+2)^2}{225}=1.
\]
En développant : $25x^2-50x+25+9y^2+36y+36=225$, soit
$\boxed{\,25x^2+9y^2-50x+36y-164=0\,}$. L'ensemble des points $A_\alpha$ est donc la conique
$(\epsilon)$ de cette équation.

\noindent\textbf{4) a)} Changement de repère $\Omega(1\,;\,-2)$ : $X=x-1$, $Y=y+2$. L'équation
$\left(\dfrac{x-1}{3}\right)^2+\left(\dfrac{y+2}{5}\right)^2=1$ devient
\[
\boxed{\,\frac{X^2}{9}+\frac{Y^2}{25}=1\,}.
\]

\noindent\textbf{4) b)} Sous la forme $\dfrac{Y^2}{25}+\dfrac{X^2}{9}=1$ (avec $25>9$), $(\epsilon)$
est une \textbf{ellipse} de centre $\Omega(1\,;\,-2)$, de grand axe porté par $\vec\jmath$. On a
$a^2=25$, $b^2=9$, donc $a=5$, $b=3$ et $c=\sqrt{a^2-b^2}=\sqrt{16}=4$. Son excentricité est
$e=\dfrac ca=\dfrac45$. Les sommets, dans le repère $(\Omega\,;\,\vec\imath,\vec\jmath)$, sont
\[
A(0\,;\,5),\quad A'(0\,;\,-5),\quad C(3\,;\,0),\quad C'(-3\,;\,0).
\]

\noindent\textbf{4) c)} Construction de l'ellipse $(\epsilon)$ de centre $\Omega$, demi-grand axe
$5$ (vertical) et demi-petit axe $3$ (horizontal).
\begin{illus}
\begin{tikzpicture}[scale=0.5]
\draw[help lines,onbline] (-4,-6) grid (4,6);
\draw[->] (-4.2,0)--(4.2,0) node[right]{$\vec\imath$};
\draw[->] (0,-6.2)--(0,6.2) node[above]{$\vec\jmath$};
\draw[onbo,very thick] (0,0) ellipse (3 and 5);
\fill[onbblue] (0,5) circle (3pt) node[above right,black]{$A$};
\fill[onbblue] (0,-5) circle (3pt) node[below right,black]{$A'$};
\fill[onbblue] (3,0) circle (3pt) node[above right,black]{$C$};
\fill[onbblue] (-3,0) circle (3pt) node[above left,black]{$C'$};
\fill[black] (0,0) circle (2.5pt) node[below left]{$\Omega$};
\node[onbo] at (-2.6,4.2) {$(\epsilon)$};
\end{tikzpicture}
\end{illus}
\fig{Ellipse $(\epsilon)$ de centre $\Omega(1\,;\,-2)$, sommets $A,A',C,C'$.}

\noindent\textbf{5) a)} L'écriture complexe d'une similitude directe est $z'=az+b$ ($a\ne0$). Le
centre $\Omega$ est invariant : $S(\Omega)=\Omega$. Avec $z_\Omega=1-2\mathrm i$, $z_C=3$,
$z_B=-\mathrm i$ :
\[
\begin{cases}1-2\mathrm i=a(1-2\mathrm i)+b\\ -\mathrm i=3a+b\end{cases}
\]
En soustrayant : $1-2\mathrm i-(-\mathrm i)=a(1-2\mathrm i)-3a$, soit
$1-\mathrm i=a(-2-2\mathrm i)=-2(1+\mathrm i)a$. Donc
\[
a=\frac{1-\mathrm i}{-2(1+\mathrm i)}=\frac{(1-\mathrm i)^2}{-2(1+\mathrm i)(1-\mathrm i)}=\frac{-2\mathrm i}{-2\cdot2}=\frac{\mathrm i}{2}.
\]
Puis $b=-\mathrm i-3a=-\mathrm i-\dfrac{3\mathrm i}{2}=-\dfrac{5\mathrm i}{2}$. D'où
\[
\boxed{\,z'=\frac{\mathrm i}{2}\,z-\frac{5\mathrm i}{2}\,}.
\]

\noindent\textbf{5) b)} L'angle de $S$ est $\arg(a)=\arg\!\left(\dfrac{\mathrm i}{2}\right)=\dfrac\pi2$.

\bigskip
\noindent\textbf{Exercice 3.}

\noindent\textbf{1)} $\overrightarrow{AB}(1\,;\,1\,;\,2)$ et $\overrightarrow{AC}(0\,;\,2\,;\,2)$.
Leurs coordonnées ne sont pas proportionnelles, donc $\overrightarrow{AB}$ et
$\overrightarrow{AC}$ ne sont pas colinéaires : $A$, $B$, $C$ ne sont pas alignés et définissent
un plan $(ABC)$.

\noindent\textbf{2)} Un vecteur normal est $\vec n=\overrightarrow{AB}\wedge\overrightarrow{AC}$.
Pour $M(x\,;\,y\,;\,z)\in(ABC)$, $\det(\overrightarrow{AB},\overrightarrow{AC},\overrightarrow{AM})=0$ :
\[
\begin{vmatrix}1&0&x+1\\ 1&2&y+1\\ 2&2&z\end{vmatrix}=0
\;\Longrightarrow\;1\bigl(2z-2(y+1)\bigr)-0+(x+1)(2-4)=0,
\]
soit $2z-2y-2-2x-2=0$, d'où $-2x-2y+2z-4=0$, c'est-à-dire
\[
\boxed{\,x+y-z+2=0\,}.
\]

\noindent\textbf{3)} Soit $(P)\,:\,x+y-z+2=0$, de vecteur normal $\vec n(1\,;\,1\,;\,-1)$, et soit
$f$ la réflexion de plan $(P)$. Pour $M(x,y,z)$ et $M'(x',y',z')=f(M)$, on a
$\overrightarrow{MM'}=\lambda\vec n$ et le milieu $I$ de $[MM']$ appartient à $(P)$. De
$\overrightarrow{MM'}=\lambda\vec n$ :
\[
x'=x+\lambda,\quad y'=y+\lambda,\quad z'=z-\lambda.\qquad(1)
\]
La condition $I\in(P)$ s'écrit $\dfrac{x+x'}{2}+\dfrac{y+y'}{2}-\dfrac{z+z'}{2}+2=0$. En y
reportant $(1)$ :
\[
\frac{2x+\lambda}{2}+\frac{2y+\lambda}{2}-\frac{2z-\lambda}{2}+2=0
\;\Longrightarrow\;3\lambda+2x+2y-2z+4=0
\;\Longrightarrow\;\lambda=-\frac23x-\frac23y+\frac23z-\frac43.
\]
En reportant dans $(1)$ :
\[
\boxed{\;
\begin{cases}
x'=\dfrac13(x-2y+2z-4)\\[4pt]
y'=\dfrac13(-2x+y+2z-4)\\[4pt]
z'=\dfrac13(2x+2y+z+4)
\end{cases}\;}
\]

\noindent\textbf{4) a)} $M(x,y,z)$ invariant par $g$ : $g(M)=M$ donne
\[
\begin{cases}3x=-x+2y-2z+4\\ 3y=2x-y-2z+4\\ 3z=-2x-2y-z+8\end{cases}
\Longleftrightarrow
\begin{cases}4x-2y+2z=4\\ -2x+4y+2z=4\\ 2x+2y+4z=8\end{cases}
\]
La troisième équation est la somme des deux premières ; le système se réduit à
$\begin{cases}2x-y+z=2\\ -x+2y+z=2\end{cases}$. En résolvant ($x=-k+2$, $y=-k+2$, $z=k$) :
\[
M\in(D)\Longleftrightarrow
\begin{cases}x=2-k\\ y=2-k\\ z=k\end{cases},\ k\in\R.
\]
C'est la droite passant par $B(0\,;\,0\,;\,2)$ (obtenu pour $k=2$) et de vecteur directeur
$\vec u(-1\,;\,-1\,;\,1)$.

\noindent\textbf{4) b) i)} Pour $g(M)=M'$, en utilisant l'expression analytique de $g$ on calcule
\[
\overrightarrow{MM'}=\begin{pmatrix}-\tfrac43x+\tfrac23y-\tfrac23z+\tfrac43\\[2pt]
\tfrac23x-\tfrac43y-\tfrac23z+\tfrac43\\[2pt]
-\tfrac23x-\tfrac23y-\tfrac43z+\tfrac83\end{pmatrix}.
\]
Alors $\overrightarrow{MM'}\cdot\vec u=$ (en multipliant les deux premières composantes par $-1$ et
la troisième par $1$) $=0$ après simplification. Donc $\overrightarrow{MM'}\perp\vec u$ :
$\overrightarrow{MM'}$ est un vecteur normal à $(D)$.

\noindent\textbf{4) b) ii)} Le milieu de $[MM']$ a pour coordonnées
$\left(\dfrac{x+x'}{2}\,;\,\dfrac{y+y'}{2}\,;\,\dfrac{z+z'}{2}\right)$. Le calcul donne
\[
\frac{x+x'}{2}=\frac13(x+y-z+2),\quad
\frac{y+y'}{2}=\frac13(x+y-z+2),\quad
\frac{z+z'}{2}=\frac13(-x-y+z+4).
\]
Ces coordonnées vérifient le système $\begin{cases}2x-y+z=2\\ -x+2y+z=2\end{cases}$ caractérisant
$(D)$ : le milieu de $[MM']$ appartient donc à $(D)$.

\noindent\textbf{4) c)} $(D)$ est l'ensemble des points invariants de $g$, et c'est la médiatrice
de $[MM']$ pour tout $M$. Donc $g$ est le \textbf{demi-tour} (retournement) d'axe $(D)$.

\noindent\textbf{5) a)} $(D)$ est dirigée par $\vec u(-1\,;\,-1\,;\,1)$, qui est colinéaire au
vecteur normal $\vec n(1\,;\,1\,;\,-1)$ de $(P)$. Donc $(D)\perp(P)$.

\noindent\textbf{5) b)} $(P)$ passe par $B$ et $(D)$ lui est perpendiculaire en $B$. La composée
d'une réflexion de plan $(P)$ et d'un demi-tour d'axe $(D)$ perpendiculaire à $(P)$ est la
symétrie centrale de centre $B=(P)\cap(D)$ :
\[
f\circ g=S_{(P)}\circ S_{(D)}=S_{(P)\cap(D)}=S_B.
\]
Donc $f\circ g$ est la \textbf{symétrie centrale de centre $B(0\,;\,0\,;\,2)$}.

\bigskip
\noindent\textbf{Partie B — Évaluation des compétences.}

\smallskip
\noindent\textbf{Tâche 1 — Masse maximale.}
\begin{remarque}
L'énoncé comporte une ambiguïté : l'unité de l'élongation n'est pas précisée et les deux données
du ressort, prises isolément, conduisent à des réponses différentes. On présente les deux
résolutions ; on travaille dans le système international (mètre, kilogramme, seconde).
\end{remarque}

\emph{Première proposition} (à partir de $mg=k\,\Delta\ell_0$ avec l'allongement
\SI{2}{\centi\metre} pour \SI{4}{\kilo\gram}). Avec $\Delta\ell_0=0{,}02$~m :
\[
4g=k\times0{,}02\;\Longrightarrow\;k=\frac{4\times9{,}5}{0{,}02}=1900\ \text{N·m}^{-1}.
\]
La masse maximale, pour l'allongement maximal $\Delta\ell_0=0{,}07$~m, est
\[
m=\frac{k\,\Delta\ell_0}{g}=\frac{1900\times0{,}07}{9{,}5}=14\ \text{kg}.
\]
On obtient $m=\SI{14}{\kilo\gram}$ (valeur plausible).

\emph{Deuxième proposition} (à partir de l'équation différentielle). L'équation
$x''(t)+\dfrac k4 x(t)=0$ a pour solutions $x(t)=a\cos\!\left(\dfrac{\sqrt k}{2}t\right)+b\sin\!\left(\dfrac{\sqrt k}{2}t\right)$. Le ressort est lâché sans vitesse initiale : $x'(0)=0\Rightarrow b=0$,
donc $x(t)=a\cos\!\left(\dfrac{\sqrt k}{2}t\right)$ avec $a\ne0$. Après $60$~s, le solide repasse pour
la première fois au point initial : $x(60)=0$, soit $\cos\!\left(30\sqrt k\right)=0$, d'où le premier
cas $30\sqrt k=\dfrac\pi2$ ($r=0$), donc $k=\left(\dfrac\pi{60}\right)^2$. Alors
\[
m=\frac{k\,\Delta\ell_0}{g}=\frac{\left(\frac{\pi}{60}\right)^2\times0{,}07}{9{,}5}=\frac{3{,}14^2\times0{,}07}{3600\times9{,}5}\approx2{,}02\times10^{-5}\ \text{kg},
\]
valeur peu plausible, ce qui confirme l'ambiguïté de l'énoncé.

\smallskip
\noindent\textbf{Tâche 2 — Chiffre d'affaires (ajustement de Mayer).}
On subdivise la série en deux sous-séries de $5$ points. Points moyens :
\[
G_1\bigl(\overline x_1\,;\,\overline y_1\bigr):\quad\overline x_1=\frac{31{,}4}{5}=6{,}82,\quad\overline y_1=\frac{1146}{5}=229{,}2,
\]
\[
G_2\bigl(\overline x_2\,;\,\overline y_2\bigr):\quad\overline x_2=\frac{52{,}8}{5}=10{,}56,\quad\overline y_2=\frac{1287}{5}=257{,}4.
\]
La droite de Mayer $y=ax+b$ passe par $G_1$ et $G_2$. Par soustraction :
\[
257{,}4-229{,}2=a(10{,}56-6{,}82)\;\Longrightarrow\;28{,}2=3{,}74\,a\;\Longrightarrow\;a\approx7{,}5401,
\]
puis $b=229{,}2-6{,}82\times7{,}5401\approx177{,}7765$. D'où $y=7{,}5401\,x+177{,}7765$.
Pour des frais de publicité de $\num{130000}$~FCFA, soit $x=13$ dizaines de milliers :
\[
y=7{,}5401\times13+177{,}7765\approx275{,}80,
\]
exprimé en dizaines de millions, soit un chiffre d'affaires estimé d'environ
$\mathbf{2{,}758\ \text{milliards de FCFA}}$.

\begin{remarque}
La correction officielle utilise $x=15$ et trouve $y=7{,}5401\times15+177{,}7765\approx290{,}88$,
soit environ $\num{2908480000}$~FCFA. On retient ici $x=13$, cohérent avec l'investissement de
$\num{130000}$~FCFA annoncé dans la situation (dépenses exprimées en dizaines de milliers).
\end{remarque}

\smallskip
\noindent\textbf{Tâche 3 — Prochaine coïncidence des fournisseurs.}
Le premier fournisseur passe tous les $21$ jours (au marché le 20/12/2020), le second tous les
$16$ jours (au marché le 27/12/2020, soit $7$ jours plus tard). En notant $x$ et $y$ les nombres de
visites menant à la prochaine coïncidence, on cherche $21x-16y=7$ avec $x,y\in\N$.

\emph{Résolution de $21x-16y=7$.} L'algorithme d'Euclide donne $21\wedge16=1$ et la relation de
Bézout $-3\times21-4\times(-16)=1$. En multipliant par $7$ :
$21\times(-21)-16\times(-28)=7$, donc $(-21\,;\,-28)$ est une solution particulière. Les solutions
entières sont
\[
S=\{(-21+16k\,;\,-28+21k)\,/\,k\in\Z\}=\{(11+16k\,;\,14+21k)\,/\,k\in\Z\}.
\]
Les plus petits entiers positifs sont $(x\,;\,y)=(11\,;\,14)$ : le premier fournisseur effectuera
$11$ visites et le second $14$. Le nombre de jours écoulés est
$21\times11=231$ jours après le 20/12/2020 (et $16\times14=224$ jours après le 27/12/2020).

En comptant $231$ jours à partir du 20 décembre 2020 (déc.~: $11$ ; +janv.~$31\to42$ ; +févr.~$28\to70$ ;
+mars~$31\to101$ ; +avr.~$30\to131$ ; +mai~$31\to162$ ; +juin~$30\to192$ ; +juil.~$31\to223$),
il reste $231-223=8$ jours en août. La prochaine coïncidence des deux fournisseurs aura donc lieu
le \textbf{8 août 2021}.
